% ------------------------------------------------------------------------------

If the \texttt{/encs} software tree does not have a required software
instantaneously available, another option is to run Singularity
containers. We run EL7 flavor of Linux, and if some projects
require Ubuntu or other distributions, there is a possibility
to run that software as a container, including the ones
translated from Docker.

The example
  \href{https://github.com/NAG-DevOps/speed-hpc/blob/master/src/lambdal-singularity.sh}
  {lambdal-singularity.sh}
showcases an immediate use of a container built for the Ubuntu-based
LambdaLabs software stack, originally built as a Docker image then
pulled in as a Singularity container that is immediately available
for use as that job example illustrates. The source material
used for the docker image was our fork of their official
repo: \url{https://github.com/NAG-DevOps/lambda-stack-dockerfiles}

NOTE: It is important if you make your own containers or pull from
DockerHub, use your \verb+/speed-scratch/$USER+ directory as these
images may easily consume gigs of space in your home directory
and you'd run out of quota there very fast.

%\noindent
TIP: To check for your quota, and the corresponding commands
to find big files, see: \url{https://www.concordia.ca/ginacody/aits/encs-data-storage.html}

We likewise built equivalent OpenISS (\xs{sect:openiss-examples})
containers from their Docker counter parts as they were used for teaching and
research~\cite{oi-containers-poster-siggraph2023}. The
images from \url{https://github.com/NAG-DevOps/openiss-dockerfiles}
and their DockerHub equivalents \url{https://hub.docker.com/u/openiss}
are found in the same public directory on \verb+/speed-scratch/nag-public+
as the LambdaLabs Singularity image. They all have \texttt{.sif} extension.
Some of them can be ran in both batch or interactive mode, some
make more sense to run interactively. They cover some basics with CUDA,
OpenGL rendering, and computer vision tasks as examples from the OpenISS
library and other libraries, including the base images that use different
distros. We also include Jupyter notebook example with Conda support.

\begin{verbatim}
/speed-scratch/nag-public:

openiss-cuda-conda-jupyter.sif
openiss-cuda-devicequery.sif
openiss-opengl-base.sif
openiss-opengl-cubes.sif
openiss-opengl-triangle.sif
openiss-reid.sif
openiss-xeyes.sif
\end{verbatim}

The currently recommended version of Singularity is
\texttt{singularity/3.10.4/default}.

This section comprises an introduction to working with Singularity, its 
containers, and what can and cannot be done with Singularity on the ENCS 
infrastructure. It is not intended to be an exhaustive presentation of 
Singularity: the program's authors do a good job of that here:
\url{https://www.sylabs.io/docs/}. It also assumes that you have successfully installed 
Singularity on a user-managed/personal system (see next paragraph as to why).

Singularity containers are essentially either built from an existing 
container, or are built from scratch. Building from scratch requires a recipe 
file (think of like a Dockerfile), and the operation \emph{must} be effected as root.
You will not have root on the ENCS infrastructure, so any built-from-scratch containers must be created 
on a user-managed/personal system. Root-level permissions are also required
(in some cases, essential; in others, for proper build functionality) for 
building from an existing container.
%, with one exception (see, Building A 
%Container From An Existing Container).
%
Three types of Singularity containers can be built: file-system; sandbox; 
squashfs. The first two are ``writable'' (meaning that changes can persist 
after the Singularity session ends). File-system containers are built around 
the ext3 file system, and are a read-write ``file'', sandbox containers are 
essentially a directory in an existing read-write space, and squashfs 
containers are a read-only compressed ``file''. Note that file-system 
containers \emph{cannot} be resized once built.

Note that the default build is a squashfs one. Also note what Singularity's 
authors have to say about the builds, ``A common workflow is to use the 
``sandbox'' mode for development of the container, and then build it as a 
default (squashfs) Singularity image when done.'' File-system containers are 
considered to be, ``legacy'', at this point in time. When built, a \emph{very small}
overhead is allotted to a file-system container (think, MB), and that 
\emph{cannot} be changed.

Probably for the most of your workflows you might find there is a
Docker container exists for your tasks, in this case you
can use the docker pull function of Singularity as a part
of you virtual environment setup as an interactive job
allocation:

\small
\begin{verbatim}
salloc --gpus=1 -n8 --mem=4Gb -t60
cd /speed-scratch/$USER/
singularity pull openiss-cuda-devicequery.sif docker://openiss/openiss-cuda-devicequery
INFO:    Converting OCI blobs to SIF format
INFO:    Starting build...
\end{verbatim}
\normalsize

\noindent
This method can be used for converting Docker containers directly on Speed.
On GPU nodes make sure to pass on the \option{--nv} flag to Singularity,
so its containers could access the GPUs. See the linked example.
